\documentclass[11pt,a4paper]{article}
\usepackage[T1]{fontenc}
\usepackage[utf8]{inputenc}
\usepackage{amsmath,amssymb,xcolor}
\usepackage[margin=25mm]{geometry}
\usepackage[colorlinks=true,urlcolor=blue]{hyperref}
\definecolor{nexorablue}{HTML}{2459B5}
\setlength{\parindent}{0pt}
\setlength{\parskip}{7pt}
\setlength{\emergencystretch}{2em}
\renewcommand{\familydefault}{\sfdefault}
\begin{document}

{\small\bfseries\color{nexorablue} NEXORA CPGE}\par
{\small MP / NEX-SLC-05 / Corrigé}\par
{\LARGE\bfseries Une perturbation ajoutée en sortie de la chaîne directe}\par
\vspace{0.25em}\hrule\vspace{0.6em}

\section*{Superposition}
On trouve $Y/R=5/(p+7)$ lorsque $D=0$, et $Y/D=(p+2)/(p+7)$ lorsque $R=0$. Les deux transferts ont un pôle à $-7$.

\section*{Régime permanent}
La contribution de la consigne vaut $5/7$. Celle de la perturbation vaut $(2/7)\times0{,}70=0{,}20$. La sortie finale est donc $5/7+0{,}20\simeq0{,}914$.


\end{document}
